\documentclass[UTF8]{ctexart}
\usepackage{xeCJK}
\usepackage{geometry}
\geometry{left=2cm,right=2cm,top=2.5cm,bottom=2.5cm}
\setlength{\headheight}{12.7pt}

\usepackage{titlesec}
\titleformat{\section}
  {\normalfont\large\bfseries}
  {\thesection}
  {1em}
  {}
\titlespacing*{\section}
  {0pt}
  {3.5ex plus 1ex minus .2ex}
  {2.3ex plus .2ex}

\usepackage{amsmath}
\usepackage{amssymb}
\usepackage{amsthm}
\usepackage{enumitem}
\setlist[enumerate]{itemsep=0pt,topsep=0pt,parsep=0pt,leftmargin=0pt,itemindent=2.5em}
\setlist[itemize]{itemsep=0pt,topsep=0pt,parsep=0pt,leftmargin=0pt,itemindent=2.5em}
\usepackage{fancyhdr}
\pagestyle{fancy}
\fancyhf{}
\usepackage{graphicx}
\usepackage{hyperref}
\usepackage{indentfirst}
\usepackage{lmodern}

\begin{document}
\setlength{\abovedisplayskip}{1pt}
\setlength{\belowdisplayskip}{1pt}

\section{预序关系}

序关系是\href{../../02 集合的一般构造/02 二元关系#nameddest=sec:定义1.1}{二元关系},
常用$\leqslant$符号表示, 把$(x,y)\in\,\,\leqslant$记为$x\leqslant y$.

此时规定它衍生的符号:
\begin{enumerate}
    \item $x\geqslant y\overset{\mathrm{def}}{\leftrightarrow}y\leqslant x$,
    \item $x<y\overset{\mathrm{def}}{\leftrightarrow}x\leqslant y\land x\neq y$,
    \item $x>y\overset{\mathrm{def}}{\leftrightarrow}y\leqslant x\land x\neq y$.
\end{enumerate}

\vspace{10pt}

最简单的序关系是预序关系.

\vspace{10pt}

\hypertarget{sec:定义1.1}{}
\noindent\textbf{定义1.1 (预序关系)}

给定集合$a$和$a$上的二元关系$\leqslant$. 如果以下两条成立:
\begin{enumerate}
    \item 自反性: $\forall x\in a:\,x\leqslant x$,
    \item 传递性: $\forall x,y,z\in a:\,
    x\leqslant y\land y\leqslant z\to x\leqslant z$,
\end{enumerate}
那么称$\leqslant$是$a$上的一个\textbf{预序关系}.
也称$(a,\leqslant)$是一个\textbf{预序结构}.

\vspace{10pt}

\hypertarget{sec:定理1.1}{}
\noindent\textbf{定理1.1}

给定预序结构$(a,\leqslant)$, 则$a$上的二元关系
\[
    \sim\,=\{(x,y)\in a\times a\mid x\leqslant y\land y\leqslant x\}
\]
是等价关系.

\noindent{\kaishu 证明.}

\begin{enumerate}
    \item 对任意的$x\in a$有$x\leqslant x$, 即$x\sim x$.
    \item 对任意的$x,y\in a$, 显然有$x\sim y\to y\sim x$.
    \item 对任意的$x,y,z\in a$, 如果$x\sim y$且$y\sim z$, 即
    \[
        x\leqslant y\land y\leqslant x\quad\land\quad
        y\leqslant z\land z\leqslant y,
    \]
    那么$x\leqslant z$且$z\leqslant x$, 即$x\sim z$. \hfill $\qedsymbol$
\end{enumerate}

\vspace{10pt}

\hypertarget{sec:定义1.2}{}
\noindent\textbf{定义1.2 (极值与最值)}

给定预序结构$(a,\leqslant)$和$m\in a$.
\begin{enumerate}
    \item 如果$\forall x\in a:x\nless m$, 就称$m$是$a$的一个\textbf{极小元},
    \item 如果$\forall x\in a:x\ngtr m$, 就称$m$是$a$的一个\textbf{极大元},
    \item 如果$\forall x\in a:x\geqslant m$, 就称$m$是$a$的\textbf{最小元},
    记作$\min a$,
    \item 如果$\forall x\in a:x\leqslant m$, 就称$m$是$a$的\textbf{最大元},
    记作$\max a$.
\end{enumerate}

\vspace{10pt}

\hypertarget{sec:定义1.3}{}
\noindent\textbf{定义1.3 (界与确界)}

给定预序结构$(a,\leqslant)$, $b\subset a$和$m\in a$.
\begin{enumerate}
    \item 如果$\forall x\in b:x\geqslant m$, 就称$m$是$b$的一个\textbf{下界},
    \item 如果$\forall x\in b:x\leqslant m$, 就称$m$是$b$的一个\textbf{上界},
    \item 如果$m$是$b$的最大下界, 就称$m$是$b$的\textbf{下确界}, 记作$\inf b$,
    \item 如果$m$是$b$的最小上界, 就称$m$是$b$的\textbf{上确界}, 记作$\sup b$.
\end{enumerate}





\newpage

\section{有向集}

\hypertarget{sec:定义2.1}{}
\noindent\textbf{定义2.1 (有向集)}

给定预序结构$(a,\leqslant)$, 如果
\[
    \forall x,y\in a\,\,\,\,\exists z\in a\,:\,x\leqslant z\land y\leqslant z,
\]
就称$(a,\leqslant)$是\textbf{有向集}, 或\textbf{有向预序集}.

\vspace{10pt}

有向集的任意两个元素都有公共上界. 进一步, 其有限子集都有上界.

\vspace{10pt}

\hypertarget{sec:定理2.1}{}
\noindent\textbf{定理2.1}

任给有向集$(a,\leqslant)$. 对$a$中任意有限个$(x_i)_{i<n}$,
$\{x_i\}_{i<n}$有上界.

\noindent{\kaishu 证明.}

首先, $n=1$的情况显然;

其次, 假设任取$(x_i)_{i<n}$, $\{x_i\}_{i<n}$有上界. 那么任取$(x_i)_{i<n+1}$,
取$\{x_i\}_{i<n}$的上界$y$, 此时存在$z$使得
\[
    y\leqslant z\quad\land\quad x_n\leqslant z,
\]
显然$z$是$\{x_i\}_{i<n+1}$的上界. \hfill $\qedsymbol$





\vspace{20pt}

\section{共尾子集}

\hypertarget{sec:定义3.1}{}
\noindent\textbf{定义3.1 (共尾子集)}

给定预序结构$(a,\leqslant)$和$s\subset a$, 如果
\[
    \forall x\in a\,\,\,\,\exists y\in s\,:\,y\geqslant x,
\]
就称$s$是$a$的\textbf{共尾子集}, 或称$s$在$a$中\textbf{共尾}.

\vspace{10pt}

\hypertarget{sec:定义3.2}{}
\noindent\textbf{定义3.2 (共尾映射)}

给定预序结构$(a,\leqslant_a),(b,\leqslant_b)$和映射$f:a\to b$.
如果$\mathrm{Ran}\,f$在$b$中共尾, 即
\[
    \forall x\in b\,\,\,\,\exists y\in a\,:\,f(y)\geqslant_b x,
\]
就称$f$是\textbf{共尾映射}.

\vspace{10pt}

显然当$s\subset a$在$a$中共尾时, 嵌入映射是共尾映射.

\end{document}
